\section{Coulomb Gases in the Complex Plane}

Under the appropriate conditions, a Coulomb Gas $\p_N$ is the Boltzmann-Gibbs probability measure given in $(\R^d)^N$. The measure $\p_N$ models an interacting gas of charged indistinguishable particles under some external field. Particles are at positions $\mmany{x}{N} \in \Se$ of dimension $d$ embedded in some space $\R^n$. We write the measure 
\begin{equation}
	\dd \p_N(\mmany{x}{N}) = \frac{e^{-\beta N^2 \Hf_N(\vec{x})}}{Z_{N,\beta}} \mcmany{\dd x}{N}{},
	\label{Equation: Coulomb Gas Measure}
\end{equation}
where $$\Hf_N(\vec{x}) = \frac{1}{N} \sum_{i = 1}^{N} \V(x) + \frac{1}{2N^2} \sum_{i \neq j}^{N} \g(x_i - x_j)$$ is usually called the Hamiltonian , or the energy of the system. The function $\V \colon \Se \mapsto \R$ is the external potential and $\g \colon \Se \mapsto (-\infty, \infty]$ the Coulomb Interacting Kernel, solution to the Poisson equation given by $- \nabla g(\vec{x}) = c_n\delta_0$. The factor $\beta N^2$ is called the inverse temperature. Take $V$ so that \eqref{Equation: Coulomb Gas Measure} is a finite measure, i.e., the normalizing constant $Z_{N, \beta}$ be finite for all $N$ and the measure support be compact.

Recall the expression for invariant ensembles and note that, for a appropriate $\V \colon \R \rightarrow \R$, we can take the appropriate dimensions for $d=1$ and $n = 2$ to recover the, for example, gaussian ensemble measure
\begin{equation}
	\p_N(\vec{x}) = \frac{e^{-\beta_N \Hf_N(\vec{x})}}{Z_{N,\beta}}, \ \ \Hf_N(\vec{x}) = \frac{1}{N} \sum_{i = 1}^{N} \V(x_i) + \frac{1}{N^2} \sum_{i < j}^{N} \log{\frac{1}{|x_i - x_j|}}.
\end{equation}
In this case, we're dealing with particles on the plane confined to the real line. However, the gases measure accepts a natural extension for any potential that we call acceptable. Together with the dimensions, we can recover other random matrix models. On the complex plane one could write explicitly
\begin{equation}
	\dd \p_N^{(\beta)} (z_1, \cdots, z_N) \deff \frac{1}{Z_{N,\beta}} \prod_{j > k = 1}^{N} |z_j - z_k|^\beta \prod_{j=1}^{N} \ee^{- \frac{\beta N}{2} Q(z_j)} \dd A(z_j)
\end{equation}
where $\dd A(z) \deff \dd^2 z / \pi$ is the area measure. This measure would imply the normalization constant
\begin{equation}
	Z_N^{(\beta)} \deff \int_{\C^N} \prod_{j > k = 1}^{N} |z_j - z_k|^\beta \prod_{j=1}^{N} \ee^{- \frac{\beta N}{2} Q(z_j)} \dd A(z_j)
\end{equation}
this is the expression, if $\beta = 2$ for the measure of the \textbf{random normal matrix}. For $Q(z) = |z|^2$ this would be the so-called \textbf{complex Ginibre ensemble}.

We could also define the measure for the configurational canonical Coulomb gas ensemble in the upper-half plane  - which has an additional complex conjugation symmetry and is governed by the law 
\begin{equation}
	\dd \tilde{\p}_N^{(\beta)} (z_1, \cdots, z_N) \deff \frac{1}{\tilde{Z}_{N,\beta}} \prod_{j > k = 1}^{N} |z_j - z_k|^\beta |z_j - \bar{z}_k|^\beta \prod_{j=1}^{N} |z_j - \bar{z}_j|^\beta  \ee^{- \frac{\beta N}{2} Q(z_j)} \dd A(z_j)
\end{equation}
and would imply the normalization constant
\begin{equation}
		\tilde{Z}_N^{(\beta)} \deff \int_{\C^N} \prod_{j > k = 1}^{N} |z_j - z_k|^\beta |z_j - \bar{z}_k|^\beta \prod_{j=1}^{N} |z_j - \bar{z}_j|^\beta  \ee^{- \frac{\beta N}{2} Q(z_j)} \dd A(z_j)
\end{equation}
and is related, for $\beta = 2$ to the \textbf{planar symplectic ensemble}. For $Q(z) = |z|^2$ this would be the so-called \textbf{symplectic Ginibre ensemble}.

For now on, we will only consider $\beta = 2$.

\section{Droplet Results}

We will use some potential $Q(z) : \C \rightarrow \R$ 
 which satisfies some nice growth condition that can be later discussed. If one takes $q(|z|) = Q(z)$, then we can say that those growth conditions are equivalent to $r q'(r)$ being increasing. Using \eqref{Proposition: Partition function}, one can write the free energy
\begin{equation}
	\ln Z_N = \ln{N!} - 2 \sum_{j=0}^{N-1} \log{h_j},  \qquad \ln \tilde{Z}_N = \ln{N!} - 2 \sum_{j=0}^{N-1} \log{2 \tilde{h}_{2j + 1}}
	\label{Equation: Partition Function}
\end{equation}
where
\begin{equation}
	h_k \deff \int_{\C} |z|^{2j} \ee^{-N Q(z)} \dd A(z), \qquad \tilde{h}_k \deff \int_{\C} |z|^{2j} \ee^{- 2 N Q(z)} \dd A(z).
\end{equation} 

\subsection{Asymptotics for random normal matrices}

The first step in the asymptotics for the normal matrices is to prepare for the Laplace method. For that we change variables such that we can identify a good function on which the method is performed. This is basically a function for which the integral values concentrates in a particular region of the integral and the rest of the domain for the integration can be controlled.

\subsubsection{Polynomial Asymptotics - Annulus}

In this case, we consider $r_0 > 0$. In order to derive the asymptotic formulas for the free energy we need to derive the asymptotics for the orthogonal polynomials and its norming constants. We first redefine our potential such that the integrand has a better expression. Using $V_{\tau(j)} = q(r) - 2 \tau(j) \log(r)$, where $\tau(j) \deff j/N \in [0, 1)$, one can rewrite the integral as
\begin{equation}
	h_j = \int_{0}^{\infty} 2r \ee^{-N V_{\tau(j)}(r)} \dd r
\end{equation}
where we have set the stage for the Laplace technique. Note that now $\tau$ will perform the function of the previous index $j$ and relates to which coefficient we're calculating. We continue as normal, separating a central part where the integral value is concentrated and another part where we can control the order of contribution. The choice of cut-size will be justified soon enough.
\begin{align}
	h_j =  \int_{|r - r_j| > \delta_N} 2r \ee^{-N V_{\tau(j)}(r)} \dd r + \int_{|r - r_j| < \delta_N} 2r \ee^{-N V_{\tau(j)}(r)} \dd r
\end{align}
As mentioned, the first integral must be somehow controlled and the second estimated as the main value for the result. Let's first show that we can control the first part of the integral, for that, consider the manipulation
\begin{align}
	\int_{|r - r_j| > \delta_N} 2r \ee^{-N V_{\tau(j)}(r)} \dd r & = \ee^{-N V_{\tau(j)}(r_{\tau(j)})} \int_{|r - r_j| > \delta_N} 2r \ee^{-N \left( V_{\tau(j)}(r) - V_{\tau(j)}(r_{\tau(j)}) \right)} \dd r \\
	& = \ee^{-N V_{\tau(j)}(r_{\tau(j)})} \int_{|r - r_j| > \delta_N} 2r \ee^{-N \left( V_{\tau(j)}(r) - V_{\tau(j)}(r_{\tau(j)}) \right)}  \ee^{-(M-M) \left( V_{\tau(j)}(r) - V_{\tau(j)}(r_{\tau(j)}) \right)} \dd r
\end{align}
Now, we need some way to estimate the difference $V_{\tau(j)}(r) - V_{\tau(j)}(r_{\tau(j)})$, for that, we first notice that $V_{\tau(j)}(r)$ has a critical (minimal) point at $r_{\tau(j)}$, the unique number which satisfies $r_\tau q(r_\tau) = 2 \tau$. As it will be relevant very soon in terms of laplace technique, it is worth mentioning the derivatives for $V_\tau(r)$,
\begin{align}
 & V'_\tau(r) = q'(r) - \frac{2\tau}{r}, \\
 & V''_\tau(r) = 4 \Delta Q(r) - \frac{1}{r} V'_\tau(r).
\end{align}
where
\begin{align}
	& V'_\tau(r(\tau)) = q'(r(\tau)) - \frac{2\tau}{r(\tau)} = 0 \\
	& V''_\tau(r(\tau)) = 4 \Delta Q(r(\tau)) - \frac{1}{r(\tau)} V'_\tau(r(\tau)) > 0.
\end{align}
Moreover, we have that
\begin{equation}
	\frac{\dd r}{\dd \tau} = \frac{2}{(r q'(r))'} |_{r = r(\tau)} = \frac{1}{2 r(\tau) \Delta Q (r(\tau))} > 0.
\end{equation}
 Another inequality needed is that we also have $r_{\tau(j+1/2)} - r_{\tau(j)} = \Boh(N^{-1})$, for some $r_{\tau(j)}< c <r_{\tau(j+1/2)}$, you can get this inequality doing the following
\begin{align}
	r_{\tau(j+1/2)} - r_{\tau(j)} & = r'(c) \left( \tau\left(j+\frac{1}{2}\right) - \tau(j)\right)	\\
	& = r'(c) \frac{j}{2N} \\
	& = \frac{j}{2N} \frac{1}{2 r(c) \Delta Q (r(c))} \\
	&  = \frac{j}{2N} \frac{2}{q'(r(c)) + r(c) q''(r(c))} \\
	& = \frac{j}{2N} \frac{2}{\frac{\dd (r q'(r))}{\dd r}|_{r = r(c)}} > 0 \\
	& = \Boh(N^{-1}).
\end{align}
Using this properties we can also deduce that $V_{\tau(j+\frac{1}{2})} - V_{\tau(j)} = \Boh(N^{-1})$, that we can show as follows
\begin{align}
	V_{\tau(j+\frac{1}{2})} - V_{\tau(j)} &= q(r_{\tau(j+\frac{1}{2})}) - q(r_{\tau(j)}) - 2 \left( \tau(j+\frac{1}{2}) \log(r_{\tau(j+\frac{1}{2})}) - \tau(j) \log(r_{\tau(j)}) \right) \\
	& = q'(c) (r_{\tau(j+\frac{1}{2})} - r_{\tau(j)}) - 2 \left( \left(r_{\tau(j+\frac{1}{2})} - r_{\tau(j)} \right) \log(r_{\tau(j+\frac{1}{2})}) + \tau(j) \log\left( \frac{r_{\tau(j+\frac{1}{2})}}{r_{\tau(j)}}  \right) \right) \\
	& = \Boh(N^{-1}).
\end{align}
And finally, because $r q'(r)$ is increasing, for all $r > r_{\tau(j)} + \delta_N$ we have
\begin{equation}
	V_{\tau(j+\frac{1}{2})}(r) \geq V_{\tau(j+\frac{1}{2})}(r_{\tau(j)} + \delta_N) =  V_{\tau(j+\frac{1}{2})}(r_{\tau(j+\frac{1}{2})}) + \frac{1}{2} V''_{\tau(j+\frac{1}{2})}(r_{\tau(j+\frac{1}{2})}) (r_{\tau(j)} + \delta_N - r_{\tau(j+\frac{1}{2})})^2 + \Boh(\delta_N^3) 
\end{equation} 
and therefore can write that
\begin{align}
	V_{\tau(j+\frac{1}{2})}(r) - V_{\tau(j)}(r_{\tau(j)}) & =  \frac{1}{2} V''_{\tau(j+\frac{1}{2})}(r_{\tau(j+\frac{1}{2})}) (r_{\tau(j)} + \delta_N - r_{\tau(j+\frac{1}{2})})^2 + \Boh(\delta_N^3) \\
	& = \delta^2 c + \Boh(N^{-1}) + \Boh(\delta_N^3) \\
	& \geq  \delta^2 c
\end{align}
for some positive real constant $c$. Getting this inequality back to the integral expression we can write that 
\begin{align}
	\int_{|r - r_j| > \delta_N} 2r \ee^{-N V_{\tau(j)}(r)} \dd r & = \ee^{-N V_{\tau(j)}(r_{\tau(j)})} \int_{|r - r_j| > \delta_N} 2r \ee^{-N \left( V_{\tau(j)}(r) - V_{\tau(j)}(r_{\tau(j)}) \right)}  \ee^{-(M-M) \left( V_{\tau(j)}(r) - V_{\tau(j)}(r_{\tau(j)}) \right)} \dd r \\
   & = \ee^{-N V_{\tau(j)}(r_{\tau(j)})}  \ee^{- c \delta_N^2 (N - M)} \int_{|r - r_j| > \delta_N} 2r \ee^{-M \left( V_{\tau(j)}(r) - V_{\tau(j)}(r_{\tau(j)}) \right)} \dd r \\
\end{align}
where you can check that
\begin{align}
	\int_{|r - r_j| > \delta_N} 2r \ee^{-M \left( V_{\tau(j)}(r) - V_{\tau(j)}(r_{\tau(j)}) \right)} \dd r & = \int_{|r - r_j| > \delta_N} 2r \ee^{-M \left( (q(r) - 2 \tau(j) \log(r)) - (q(r_{\tau(j)}) - 2 \tau(j) \log(r_{\tau(j)})) \right)} \dd r  \\
	& = 2 \left( \ee^{-q(r_{\tau(j)})} r_{\tau(j)}^{-2 \tau(j)} \right)^M \int_{|r - r_j| > \delta_N} \left( \ee^{-q(r)} r^{2 \tau(j)} \right)^M \dd r
\end{align}
where we choose $M$ to be a big enough number such that the integral converges. If that is true, 
\begin{equation}
	\int_{|r - r_j| > \delta_N} 2r \ee^{-M \left( V_{\tau(j)}(r) - V_{\tau(j)}(r_{\tau(j)}) \right)} \dd r = \ee^{-N V_{\tau(j)}(r_{\tau(j)})} \epsilon_N
\end{equation}
with $\epsilon_N = \Boh(\ee^{-c(\log(N))^2})$. Now we must deal with the second integral. The idea here is quite direct, to use the Taylor expansion for $V_{\tau(j)}(r)$. That is,  
\begin{align}
  (2) & = \ee^{-N V_{\tau(j)}(r_{\tau(j)})} \int_{|r - r_j| < \delta_N} 2r \ee^{-N \left( V_{\tau(j)}(r) - V_{\tau(j)}(r_{\tau(j)}) \right)} \dd r \\
	& =  \ee^{-N V_{\tau(j)}(r_{\tau(j)})} \int_{\delta_N}^{\delta_N} 2(r_{\tau(j)} + r') \ee^{-N \left( \left(V_{\tau(j)}(r_{\tau(j)}) + \frac{1}{2}(4 \Delta Q(r_{\tau(j)})) (r')^2 + \frac{1}{3!} V^{(3)}_{\tau(j)}(r_{\tau(j)}) (r')^3 + \frac{1}{4!} V^{(4)}_{\tau(j)}(r_{\tau(j)}) (r')^4 + \Boh(r'^5) \right) - V_{\tau(j)}(r_{\tau(j)}) \right)} \dd r' \\
	&  = \frac{2 r_{\tau(j)}}{\sqrt{N}} \ee^{-N V_{\tau(j)}(r_{\tau(j)})} \int_{\sqrt{N} \delta_N}^{\sqrt{N}\delta_N}\left(1 + \frac{\bar{r}}{r_{\tau(j)} \sqrt{N}}\right) \ee^{-2\Delta Q(r_{\tau(j)}) \bar{r}^2} \ee^{\frac{1}{3! \sqrt{N}} V^{(3)}_{\tau(j)}(r_{\tau(j)}) (\bar{r})^3 + \frac{1}{4! N} V^{(4)}_{\tau(j)}(r_{\tau(j)}) (\bar{r})^4 + \Boh(N^{3/2} \bar{r}^5)} \dd \bar{r}
\end{align}
where $r' = r_\tau(j) - r$ and $\bar{r} =  \sqrt{N}r'$. Now we take the Taylor Series for 
\begin{equation}
	\ee^{\frac{1}{3! \sqrt{N}} V^{(3)}_{\tau(j)}(r_{\tau(j)}) (\bar{r})^3 + \frac{1}{4! N} V^{(4)}_{\tau(j)}(r_{\tau(j)}) (\bar{r})^4}
\end{equation}
to get the asymptotic expansion (since odd terms vanish)
\begin{equation}
	(2) = \frac{2 r_{\tau(j)}}{\sqrt{N}} \ee^{-N V_{\tau(j)}(r_{\tau(j)})} \int_{-\infty}^{\infty} \ee^{-2\Delta Q(r_{\tau(j)}) \bar{r}^2} \left[ 1- \frac{1}{N} \left( \frac{1}{6 r_{\tau(j)}} V^{(3)}_{\tau(j)}(r_{\tau(j)}) + \frac{1}{24} V^{(4)}_{\tau(j)}(r_{\tau(j)}) \right) \bar{r}^4 + \frac{1}{N} \frac{1}{72} \left(V^{(3)}_{\tau(j)}(r_{\tau(j)}) \right)^2 \bar{r}^6 \right] \dd r + \Boh(N^{-2})
\end{equation}
the result can be directly evaluated using the Gaussian integrals
\begin{equation}
	\int_{\R} \ee^{-2ar^2} \dd r = \sqrt{\frac{\pi}{2}} \frac{1}{a^{1/2}}
\end{equation}
\begin{equation}
	\int_{\R} \ee^{-2ar^2} r^4 \dd r = \sqrt{\frac{\pi}{2}} \frac{3}{16} \frac{1}{a^{5/2}}
\end{equation}
\begin{equation}
	\int_{\R} \ee^{-2ar^2} r^6 \dd r = \sqrt{\frac{\pi}{2}} \frac{15}{64} \frac{1}{a^{7/2}}
\end{equation}

So now we need to return to  \ref{Equation: Partition Function} and calculate 
\begin{align}
	\sum_{j=0}^{N-1} \log{h_j} & = \sum_{j=0}^{N-1} -N V_{\tau(j)}(r_{\tau(j)}) + \frac{1}{2} \left(\log(2\pi r^2_{\tau(j)}) - \log N - \log \Delta Q(r_{\tau(j)}) \right) + \frac{1}{N} \mathfrak{B}_1(r_{\tau(j)}) + \Boh(N^{-2}) 
\end{align}

We first look at the first term using the Euler-Maclaurin formula
\begin{equation}
	\sum_{j=0}^{N-1} V_{\tau(j)}(r_{\tau(j)}) = \int_0^N V_{\tau(j)}(r_{\tau(j)})  \dd t - \frac{1}{2} \left( V_{\tau(j)}(r_{\tau(j)})  - V_{\tau(0)}(r_{\tau(0)}) \right) + \frac{1}{12} \left[ \partial_t V_{\tau(j)}(r_{\tau(j)}) |_{t=N} - \partial_t V_{\tau(j)}(r_{\tau(j)}) |_{t=0} \right] + \Boh(N^{-3})
\end{equation}
where the first term is solved doing a integral by parts. Setting $s = r_{\tau(j)}$, we note that $V_{\tau(j)}(r_{\tau(j)}) = q(r) - 2 \tau \log r$, with $\tau(j) = j/N$, the expression for the derivative of $r$, 
\begin{equation}
	\frac{\dd r}{\dd \tau} = \frac{1}{2 r_{\tau(j)} \Delta Q(r_{\tau(j)})} = \frac{2}{(r q'(r))'}
\end{equation}
and $r_{\tau(j)} q'(r_{\tau(j)}) = 2 \tau$. We then are able to rewrite
\begin{align}
	\int_0^N V_{\tau(j)}(r_{\tau(j)})  \dd t & = \int_0^N (q(r_{\tau(j)}) - 2 \tau(j) \log(r_{\tau(j)})) \dd t \\
	& = \int_{r_0}^{r_1} (q(s) - s q'(s) \log(s)) s \Delta Q(s) 2 N \dd s \\
	& = 2 N \left[ \int_{r_0}^{r_1} q(s) s \Delta Q(s)\dd s - \int_{r_0}^{r_1} s^2 q'(s) \log(s) \Delta Q(s) \dd s \right]\\
	& =: N \left[ \int_S Q \Delta Q \dd A - 2 \int_{r_0}^{r_1} s^2 q'(s) \log(s) \Delta Q(s) \dd s \right]
\end{align}
and finally using that $r_1 q'(r_1) = 2$ and $r_0 q'(r_0) = 0$ we can apply the integration by parts on the second term and get
\begin{align}
	2 \int_{r_0}^{r_1} s^2 q'(s) \log(s) \Delta Q(s) \dd s & = 2 \int_{r_0}^{r_1} s^2 q'(s) \log(s) \frac{1}{4s} (s q'(s))' \dd s \\
	&= \frac{1}{2} \int_{r_0}^{r_1} s q'(s) \log(s) (s q'(s))' \dd s \\
	& = \frac{1}{4} \int_{r_0}^{r_1} ((sq'(s))^2)' \log(s) \dd s \\
	& = \frac{1}{4} (s q'(s))^2 \log(s) |_{r_0}^{r_1} + \frac{1}{4}  \int_{r_0}^{r_1} (s q'(s))^2 \frac{1}{s} \dd s \\
	& = \log(r_1) + \frac{1}{4} \int_{r_0}^{r_1} s q'(s)^2 \dd s \\
	& = \log(r_1) + \frac{1}{2} q(r_1) - \frac{1}{4} \int_{r_0}^{r_1} q(s) (s q'(s))' \dd s
\end{align}
Leaving us with
\begin{align}
	\int_0^N V_{\tau(j)}(r_{\tau(j)})  \dd t & = N \left[\int_S Q \Delta Q \dd A - \log(r_1) + \frac{1}{2} q(r_1) - \frac{1}{4} \int_{r_0}^{r_1} q(s) (s q'(s))' \dd s \right]  \\
	& =: N \left[\frac{1}{2} \int_S Q \Delta Q \dd A - \log(r_1) + \frac{1}{2} q(r_1) \right]  \\
	& =: N I_Q[\mu_Q] 
\end{align}
We also need an expression for $ V_{\tau(j)}(r_{\tau(j)})  - V_{\tau(0)}(r_{\tau(0)})$. This is direct noting that
\begin{equation}
	V_{\tau(j)}(r_{\tau(j)})  - V_{\tau(0)}(r_{\tau(0)}) = V_1(r_1) - V_0(r_0) = q(r_1) - q(r_0)
 - 2 \log(r_1) =: 2 U_{\mu_Q}(0)
 \end{equation}
 And finally we can the remaining term , we use the Leibniz rule and obtain
\begin{equation}
	\partial_t V_{\tau(t)}(r_{\tau(t)})|_{t=N} = \frac{1}{N} \left[ \partial_t(q(r_{\tau(t)}) - 2 \log(r_{\tau(t)}))_{\tau=1} - 2 \log(r_1) \right]
\end{equation}
\begin{equation}
	\partial_t V_{\tau(t)}(r_{\tau(t)})|_{t=0} = \frac{1}{N} \left[ \partial_t(q(r_{\tau(t)}) - 2 \log(r_{\tau(t)}))_{\tau= 0} - 2 \log(r_0) \right]
\end{equation}
and since we can calculate
\begin{equation}
	\partial_\tau q(r_{\tau}) |_{\tau = 0} = \partial_{r_\tau} (q(r_\tau)) \frac{\dd r_\tau}{\dd \tau} = q'(r_\tau) \frac{1}{2 r_{\tau(j)} \Delta Q(r_\tau(j))} = \frac{q'(r_0)}{2 r_0 \Delta Q(r_0)}
\end{equation}
\begin{equation}
	\partial_\tau q(r_{\tau}) |_{\tau = 1} = \frac{q'(r_1)}{2 r_1 \Delta Q(r_1)}
\end{equation}
\begin{equation}
	\partial_\tau (2 \log(r_\tau)) |_{\tau = 1} = \frac{2}{r_\tau} \frac{1}{2 r_\tau \Delta Q (r_\tau)} = \frac{1}{r_1^2 \Delta Q(r_1)}
\end{equation}
We can also express the difference as
\begin{equation}
	\partial_t V_{\tau(t)}(r_{\tau(t)})|_{t=N}  - \partial_t V_{\tau(t)}(r_{\tau(t)})|_{t=0} = \frac{1}{N} \left[ 2(\log(r_0) - \log(r_1)) +  \frac{q'(r_0)}{2 r_0 \Delta Q(r_0)} +  \frac{q'(r_1)}{2 r_1 \Delta Q(r_1)} +  \frac{1}{r_1^2 \Delta Q(r_1)} \right] = \frac{2}{N} \log\left(\frac{r_0}{r_1}\right)
\end{equation}
And the proof for the summation $\sum_{j=0}^{N-1} V_{\tau(j)}(r_{\tau(j)})$ is complete if one gathers the terms. We move on to the next term from where we need to determine both $ \sum_{j=0}^{N-1}\left(\log(2\pi r^2_{\tau(j)}) - \log \Delta Q(r_{\tau(j)}) \right)$ Lets star with the second summation, to which we apply Euler-Maclaurin again
\begin{equation}
	\sum_{j=0}^{N-1}  \log \Delta Q(r_{\tau(j)}) = \int_{0}^{N} \log \Delta Q(r_{\tau(j)}) \dd t - \frac{1}{2} \log \left( \frac{\Delta Q(r_1)}{\Delta Q(r_0)} \right) + \frac{1}{12} \left[ \partial_t \log \Delta Q (r_\tau) |_{t=1} - \partial_t \log \Delta Q(r_\tau) |_{t=0}\right]
\end{equation}
Now we tackle the first integral on the right hand side, changing variables as $s = r_{\tau(t)}$ we would have
\begin{align}
	\int_{0}^{N} \log \Delta Q(r_{\tau(j)}) \dd t & = N \int_{r_0}^{r_1} \log(\Delta Q(s)) 2 s \Delta Q(s) \dd s \\
	& = N \int_{S} \log(\Delta Q) \dd \mu_Q
\end{align}
and for the third term on the right hand side, 
\begin{equation}
	\partial_t \log \Delta Q (r_\tau) |_{t=1} - \partial_t \log \Delta Q(r_\tau) |_{t=0} = \frac{1}{2N} \left( \frac{(\partial_r \Delta Q)(r_1)}{r_1 (\Delta Q(r_1))^2} - \frac{(\partial_r \Delta Q)(r_0)}{r_0 (\Delta Q(r_0))^2}\right) = \Boh(N^{-1})
\end{equation}

With all of this, one can now finally express
\begin{equation}
	\sum_{j=0}^{N-1} \log(h_j) = - N^2 I_Q[\mu_Q] - \frac{N}{2} \log N + N \left(\frac{\log(2\pi)}{2} - \frac{1}{2} E_Q[\mu_Q]\right) - \frac{1}{3} \log\left(\frac{r_1}{r_0} \right) + \frac{1}{4} \log\left(\frac{\Delta Q(r_1)}{\Delta Q(r_0)} \right) + \int_S \mathfrak{B}_1 \dd \mu_Q + \Boh(N^{-1})
\end{equation}
from which we can get the asymptotic expansion using that
\begin{equation}
	\log(N!) = N\log(N) - N + \frac{1}{2} \log(N) + \frac{1}{2} \log(2 \pi) + \Boh(N^{-1})
\end{equation}

\subsubsection{Polynomial Asymptotics - Disk}

In this case, we consider $r_0 = 0$. In the asymptotic expansions of $h_j$, one should distinguish the following two cases depending on a small constant $\epsilon > 0$. 
\begin{itemize}
	\item Case 1: $r_{j/N} >> N^{-\epsilon}$. For the annular droplet case where $r_0 > 0$, this case covers all $j = 0, 1, \cdots , N-1$. On the other hand, for the disc droplet case where $r_0 = 0$, this case covers only $j = m_N, m_N + 1, \cdots , N-1$ for $m_N = N^\epsilon$ with some $\epsilon > 0$;
	\item $r_{j/N} << N^{-\epsilon}$. This covers the remaining disc droplet case with $j = 0, 1, \cdots, m_N-1$. Notably, the asymptotic expansion involves gamma functions in this case. 
\end{itemize}

We separate this in two proofs, the first one for $j = m_N, m_N + 1, \cdots, N-1$  is very similar to our previous proof. To actually expand on the details, we need some useful facts. \textcolor{red}{First, recall tat $r_{\tau}$ satisfies $r_\tau = \Boh(\tau^{\frac{1}{2}})$ as $\tau \rightarrow 0$}. Note also that for $d \geq 1$, we can choose $C_1$ a constant uniformly for all $\tau$ and write
\begin{equation}
	|V^{(d)}_\tau(r_\tau)| = |q^{(d)}(r_\tau) + (-1)^{d} 2 \tau (d-1)! r_\tau^{-d} | \leq C_1 (1 + \tau^{-\frac{d}{2} + 1}).
\end{equation} 
We then star as usual dividing the integral into two parts
\begin{align}
	h_j & = \int_{0}^{\infty} 2 r^{2j + 1} \ee^{-N q(r)} \dd r = \int_{0}^{\infty} 2 r \ee^{-N V_{\tau(j)} (r)} \dd r \\
	 &= \int_{r_{\tau(j)} - \delta_N}^{r_{\tau(j)} + \delta_N} 2 r \ee^{-N V_{\tau(j)}(r)} \dd r+ \int_{|r - r_{\tau(j)} > \delta_N} 2 r \ee^{-NV_{\tau(j)}(r)} \dd r
\end{align}
and now we have two integrals to solve. Let's star with the outer integral, which we should be able to control the growth. For $m_N \leq j < N$, we have
\begin{equation}
	r_{\tau(j + 1/2)} - r_{\tau(j)} = \Boh((jN)^{-\frac{1}{2}}).
\end{equation}
Since the potential is increasing in $(r_\tau, \infty)$, the Taylor series expansion for $V_\tau$ is for all $r > r_{\tau(j)} + \delta_N$
\begin{align}
	V_{\tau(j + 1/2)}(r) & \geq V_{\tau(j + 1/2)}(r) \\
	& = V_{\tau(j + 1/2)}(r_{\tau(j + 1/2)}) + \frac{1}{2} V^{(2)}_{\tau(j + 1/2)}(r_{\tau(j + 1/2)}) (r_{\tau(j)}) + \delta_N - r_{\tau(j + 1/2)})^2 + \Boh(\tau(j)^{-1/2} \delta^3_N).
\end{align}
where $\Boh(\tau(j)^{-1/2} \delta^3_N) = \Boh(N^{-1} j^{-1/2}(\log(N))^3)$. Similarly, since $V_{\tau}$ is decreasing in $(0, r_\tau)$, we have for all $r < r_{\tau(j)} - \delta_N$
\begin{align}
	V_{\tau(j + 1/2)}(r) & \geq V_{\tau(j - 1/2)}(r) \\
	& = V_{\tau(j + 1/2)}(r_{\tau(j + 1/2)}) + \frac{1}{2} V^{(2)}_{\tau(j + 1/2)}(r_{\tau(j + 1/2)}) (r_{\tau(j)}) - \delta_N - r_{\tau(j + 1/2)})^2 + \Boh(\tau(j)^{-1/2} \delta^3_N).
\end{align}
Using the Taylor series for $V_\tau$ again, we have
\begin{align}
	V_{\tau(j+1/2)}(r_{\tau(j+1/2)}) - V_{\tau(j)}(r_{\tau(j)}) & = V_{\tau(j)}(r_{\tau(j+1/2)}) - V_{\tau(j)}(r_{\tau(j)}) - \frac{1}{N} \log r_{\tau(j+1/2)} \\
	& = \Boh((jN)^{-1})  - \frac{1}{N} \log r_{\tau(j+1/2)}
\end{align}
Thus, if follows from the last two inequalities that
\begin{align}
	\ee^{N V_{\tau(j)}(r_{\tau(j)})} \int_{|r - r_{\tau(j)}| > \delta_N} 2 r^{2j+1} \ee^{-N q(r)} \dd r & = \int_{|r - r_{\tau(j)}| > \delta_N} 2 \ee^{-N \left( V_{\tau(j + 1/2)}(r) - V_{\tau(j)}(r_{\tau(j)}) \right)} \dd r \\
	& \leq \ee^{-c_1 (\log(N))^2} \ee^{\frac{N-M}{N} \log(r_{\tau(j+1/2)})}  \int_{|r - r_{\tau(j)}| > \delta_N} 2 \ee^{-M \left( V_{\tau(j + 1/2)}(r) - V_{\tau(j)}(r_{\tau(j)}) \right)} \dd r = \Boh(\ee^{c_2 (\log(N))^2})
\end{align}
for some constants $c_1, c_2 > 0$. We now need to examine the integral near the critical point. That is, 
\begin{align}
	& \int_{r_{\tau(j)} - \delta_N}^{r_{\tau(j)} + \delta_N} 2r \ee^{-N V_{\tau(j)}} \dd r \\
	& =  \ee^{-N V_{\tau(j)}(r_{\tau(j)})} \int_{-\delta_N}^{\delta_N} 2(r_{\tau(j)} + t) \ee^{-N \left( \frac{1}{2} V^{(2)}_{\tau(j)}(r_{\tau(j)}) (t)^2 + \frac{1}{3!} V^{(3)}_{\tau(j)}(r_{\tau(j)}) (t)^3 + \frac{1}{4!} V^{(4)}_{\tau(j)}(r_{\tau(j)}) (t)^4 + \Boh(\tau(j)^{-\frac{3}{2}}  t^5) \right)} \dd t \\
	&  =  \ee^{-N V_{\tau(j)}(r_{\tau(j)})}  \int_{-\sqrt{N} \delta_N}^{\sqrt{N} \delta_N} 2 \ee^{-2\Delta Q(r_{\tau(j)}) t^2} (r_{\tau(j)} + \frac{t}{\sqrt{N}})  \left(  1 - \frac{V^{(3)}_{\tau(j)}(r_{\tau(j)}) (t)^3 }{6 \sqrt{N}} - \frac{V^{(4)}_{\tau(j)}(r_{\tau(j)}) (t)^4}{24 N} + \frac{V^{(3)}_{\tau(j)}(r_{\tau(j)})^2 (t)^6}{72 N} + \epsilon'_{N,1} \right) \dd t
\end{align}
where  $\epsilon_{N,1} = \Boh(j^{-\frac{3}{2}} (\log N)^\alpha)$ for some $\alpha > 0$.

Now, for $j \leq m_N$, we need to define some $r^*_\tau \deff r_\tau \log(N)$ from which we separate the integral in a important part close to the origin and another part which has a lesser order. Consider
\begin{equation}
h_j = \int_{0}^{\infty} 2r^{2j+1} \ee^{-N q(r)} \dd r =\int_{0}^{r^*_{\tau(j+\frac{1}{2})}} 2r^{2j+1} \ee^{-N q(r)} \dd r + \int_{r^*_{\tau(j+\frac{1}{2})}}^{\infty} 2r^{2j+1} \ee^{-N q(r)} \dd r
\end{equation}
\textcolor{red}{Due to strict subharmonicity of Q in a neighborhood of the droplet, $r_\tau$ satisfies $r_\tau = \Boh (\tau^{\frac{1}{2}})$ as $\tau \rightarrow 0$}. Since $V_\tau$ has a global minimum at $r = r_\tau$ and increases for larger values, for all $r > r^*_{\tau(j + 1/2)}$ we have
\begin{align}
	V_{\tau(j+1/2)}(r) & \geq V_{\tau(j+1/2)}(r)(r^*_{\tau(j+1/2)}) = q(r^*_{\tau(j+1/2)}) - \frac{2j + 1}{N} \log(r^*_{\tau(j+1/2)}) \\
	& = q(0) + \frac{1}{2} q^{(2)}(0)(r^*_{\tau(j+1/2)})^2 - \frac{2j + 1}{N} \log(r^*_{\tau(j+1/2)}) + \Boh(r^*_{\tau(j+1/2)})^3.
\end{align}
Remember that $q'(0) = 0$ and $\Delta Q(z) \in (0, \infty)$ near the origin. Using this inequality,
\begin{align}
	& \ee^{N q(0)} \int_{r^*_{\tau(j+\frac{1}{2})}}^{\infty} 2r^{2j+1} \ee^{-N q(r)} \dd r = \int_{r^*_{\tau(j+\frac{1}{2})}}^{\infty} 2 \ee^{-N (V_{\tau(j+1/2)}(r) q(0))} \dd r \\
	& \leq \ee^{-c_1 (N-M) (r^*_{\tau(j+\frac{1}{2})})^2} (r^*_{\tau(j+\frac{1}{2})})^{(2j+1) \frac{N-M}{N}} \int_{r^*_{\tau(j+\frac{1}{2})}}^{\infty} 2 \ee^{-M (V_{\tau(j+1/2)}(r) q(0))} \dd r \deff \epsilon_N(j)
\end{align}
for some $c_1 > 0$ and $M$ as given in the other proof, such that the appropriate integral converges. Since $r_\tau = \Boh(\tau^{\frac{1}{2}})$, there is a constant $c_2 > 0$ such that
\begin{equation}
	\epsilon_N(j) = \Boh\left( (r^*_{\tau(j+\frac{1}{2})})^{(2j+1) \frac{N-M}{N}} \ee^{-c_2 (j + \frac{1}{2}) (\log N)^2} \right) = \Boh\left( (\frac{2j+1}{2N} (\log(N))^2)^{(2j+1) \frac{N-M}{N}} \ee^{-c_2 (j + \frac{1}{2}) (\log N)^2} \right)
\end{equation}
Therefore we obtain
\begin{align}
	h_j & = \int_{0}^{\infty} 2r^{2j+1} \ee^{-N q(r)} \dd r = \int_{0}^{r^*_{\tau(j+\frac{1}{2})}} 2r^{2j+1} \ee^{-N q(r)} \dd + \ee^{-N q(0)} \epsilon_N(j) \\
	& = \int_{0}^{r^*_{\tau(j+\frac{1}{2})}} 2r^{2j+1} \ee^{-N \left( q(0) + \frac{1}{2} q^{(2)}(0) r^2 \right) + \Boh(N(r^*_{\tau(j+\frac{1}{2})})^3)} \dd  r + \ee^{-N q(0)} \epsilon_N(j) \\
	& =  \ee^{-N q(0)} \left[ N{-(j+1)} \int_{0}^{r^*_{\tau(j+\frac{1}{2})}} 2r^{2j+1} \ee^{-\frac{1}{2} q^{(2)}(0) r^2} \dd r \left( 1 + \Boh(N(r^*_{\tau(j+\frac{1}{2})})^3) \right)+ \epsilon_N(j) \right] \\
	&=  \ee^{-N q(0)} \left[ N{-(j+1)} \Gamma(j+1)(\frac{1}{2} q^{(2)}(0))^{-(j+1)} \left( 1 + \Boh(N^{-\frac{1}{2}} (j+\frac{1}{2})^{\frac{3}{2}} (\log N)^3) \right)+ \epsilon_N(j) \right]
\end{align}
and this completes this part of the proof.  

Now we complete the case calculating, using the results above,
\begin{equation}
	\sum_{j=0}^{N} \log(h_j) = \sum_{j=0}^{m_N - 1} \log(h_j)  + \sum_{j=m_N}^{N - 1} \log(h_j) 
\end{equation}
The first term in the right is easily determined, we need only to write
\begin{align}
	 \sum_{j=0}^{m_N - 1} \log(h_j) & = -m_N N q(0) - \frac{m_N(m_N + 1)}{2} \left( \log(N) + \log\left( \frac{1}{2} q^{(2)}(0) \right) \right) + \log(G(m_N + 1)) + \Boh(N^{-\frac{1}{2}(1-5\epsilon)} (\log(N))^3) \\
	 & = -m_N N q(0) - \frac{m_N(m_N + 1)}{2} \left( \log(N) + \log\left( \frac{1}{2} q^{(2)}(0) \right)  \right) +  \frac{m_N^2 \log(m_N)}{2} - \frac{3}{4} m_N^2 + \frac{\log(2\pi) m_N}{2} - \frac{\log(m_N)}{12} + \gamma'(-1)+ \Boh(m_N^{-2}+ N^{-\frac{1}{2}(1-5\epsilon)} (\log(N))^3)
\end{align}
where we have used the asymptotic for the G-Barnes function
\begin{equation}
	\log(G(z+1)) = \frac{z^2 \log(z)}{2} - \frac{3}{4} z^2 + \frac{\log(2\pi) z}{2} - \frac{\log(z)}{12} + \gamma'(-1) + \Boh(\frac{1}{z^2})
\end{equation}

The second summation is a bit more detailed, we use the previously defined
\begin{equation}
	\log(h_j) = -N V_{\tau(j)} (r_{\tau(j)}) + \frac{1}{2} \log{\left( \frac{2\pi r^2_{\tau(j)}}{N \Delta Q(r_{\tau(j)})} \right)} + \frac{1}{N} \mathfrak{B}_1(r_{\tau(j)}) + \Boh(j^{-\frac{3}{2}} (\log(N))^\alpha)
\end{equation}
and need to determine three things in order to get the summation,
\begin{align}
	& \sum_{j=m_N}^{N-1} V_{\tau(j)}(r_{\tau(j)}), \\
	& \sum_{j=m_N}^{N-1} \log \Delta Q(r_{\tau(j)}), \\
	&  \sum_{j=m_N}^{N-1} \log(r_{\tau(j)}).
\end{align}

For the first we proceed as before and perform the Euler-Maclaurin formula.
\begin{equation}
	\sum_{j=m_N}^{N-1} V_{\tau(j)}(r_{\tau(j)}) = \int_{m_N}^N V_{\tau(j)}(r_{\tau(j)})  \dd t - \frac{1}{2} \left( V_{\tau(N)}(r_{\tau(N)})  - V_{\tau(m_N)}(r_{\tau(m_N)}) \right) + \frac{1}{12} \left[ \partial_t V_{\tau(j)}(r_{\tau(j)}) |_{t=N} - \partial_t V_{\tau(j)}(r_{\tau(j)}) |_{t=m_N} \right] + \Boh(N^{-1-2\epsilon})
\end{equation}
where we used $\partial^3_t(V_{\tau(t)}(r_{\tau(t)})) |_{t=m_N} = \Boh(N^{-3} (\tau(m_N))^{-2}) = \Boh(m_N^{-2} N^{-1})$. The first term can be solved in a very similar way as before.
\begin{align}
	\int_{m_N}^{N} V_{\tau(j)}(r_{\tau(j)}) & = 2N \int_{r_\tau(m_N)}^{r_1} (q(s) - s q'(s) \log s) s \Delta Q(s) \dd s \\
	& = N \left( \frac{1}{2} \int_{S \backslash S_{\tau(m_N)}} Q \Delta Q \dd A - \log(r_1) + (\tau(m_N))^2 \log(r_{\tau(m_N)} + \frac{1}{2} \left( q(r_1) - \tau(m_N) q(r_{\tau(m_N)}) \right) \right) \\
	& = N I_Q[\mu_Q] - \frac{N}{2} \int_{S_{\tau(m_N)}} Q \Delta Q \dd A + \frac{m_N^2}{N} \log(r_{\tau(m_N)}) - \frac{m_N}{2} q(r_{\tau(m_N)}) 
\end{align}
where, \textcolor{red}{because 
\begin{equation}
	r_\tau = \left(\frac{2 \tau}{q^{(2)}(0)} \right)^{\frac{1}{2}} + \Boh(\tau) = \left(\frac{\tau}{\Delta Q(0)} \right)^{\frac{1}{2}} + \Boh(\tau), \qquad \text{as} \quad \tau \rightarrow 0
\end{equation}}
we have that
\begin{equation}
	\log(r_{\tau(m_N)}) = \frac{1}{2} \log{\left(\frac{\tau(m_N)}{\Delta Q(0)}\right)} + \Boh(\tau(m_N)^{\frac{1}{2}}) = \frac{1}{2} \log{\left(\frac{\tau(m_N)}{\Delta Q(0)}\right)} + \Boh(N^{-\frac{1}{2}(1-\epsilon)}) 
\end{equation}
and
\begin{equation}
	q(r_{\tau(m_N)}) = q(0) + \frac{1}{2} q^{(2)}(0)(r_{\tau(m_N)})^2 + \Boh((\tau(m_N)^{\frac{3}{2}})) = q(0) + \tau(m_N) + \Boh(N^{-\frac{3}{2}(1 - \epsilon)})
\end{equation}